\chapter{Introduction}
\label{cap:cap1}

Running services reliably across multiple machines is a common challenge in modern software systems. As user traffic grows, a single server can no longer handle all requests. The application must be split across a cluster of machines, where each machine is called a node.

Managing a cluster brings several problems. The system needs to know which nodes are available. When a new service is deployed, it must decide which node should run it. If a node crashes, the services running on it must be moved to another node. The system also needs access control and a record of who did what and when.

Kubernetes is the most popular platform for solving these problems, but it has hundreds of components and requires a team of engineers to run. This thesis proposes a simpler alternative: a distributed platform built from scratch that implements the fundamental concepts (heartbeats, scheduling, load balancing, fault tolerance, security and audit) in a clear and understandable way.

\section{Objectives}

The main objective of this thesis is to design and implement a distributed platform for running containerized services within a cluster. The platform addresses the following specific objectives:

\begin{itemize}
    \item \textbf{Automatic node discovery}. Worker nodes are discovered automatically by the master when they send their first heartbeat via UDP multicast. No explicit join or leave mechanism is needed.
    \item \textbf{Resource pools}. Nodes are grouped into pools based on their hardware resources (CPU, RAM). Services can be deployed to a specific pool, ensuring they run on appropriate hardware.
    \item \textbf{Load balancing}. When deploying a service, the platform selects the best node using a pluggable load balancing strategy (Least Connections or Weighted).
    \item \textbf{Fault tolerance}. The platform detects node failures through heartbeat monitoring and automatically restarts affected containers on healthy nodes. A Raft consensus protocol ensures master node availability.
    \item \textbf{Security}. User authentication is handled through JWT tokens and a role-based access control system (Admin, Writer, Reader) restricts what each user can do.
    \item \textbf{Audit}. All important actions are logged in an append-only audit log that can be queried through the API.
\end{itemize}

\section{Methodology and Tools}

The platform is implemented in Go, chosen for its built-in concurrency support (goroutines), efficient networking and suitability for systems programming. The main tools and libraries used are:

\begin{itemize}
    \item \textbf{Go}. The programming language for both master and worker components
    \item \textbf{Docker SDK}. For managing containers on worker nodes
    \item \textbf{HashiCorp Raft}. For leader election and state replication between master nodes
    \item \textbf{UDP multicast}. For heartbeat communication between workers and master
    \item \textbf{JWT (JSON Web Tokens)}. For user authentication
\end{itemize}

\section{Thesis Structure}

The thesis is organized into five chapters:

\begin{itemize}
    \item \textbf{Chapter 1. Introduction}: presents the context, motivation, objectives, methodology and an overview of the thesis structure.
    \item \textbf{Chapter 2. Theoretical Background and Similar Solutions}: covers the key concepts of distributed systems (heartbeats, consensus, fault tolerance, load balancing, containerization) and compares existing platforms (Kubernetes, Docker Swarm and HashiCorp Nomad) with the proposed solution. It also defines the expected characteristics of the platform.
    \item \textbf{Chapter 3. Proposed Solution}: describes the architecture of the platform in detail, including the master-worker component design, domain model, heartbeat-based node discovery, fault tolerance mechanisms, load balancing strategies, REST API, service deployment flow, JWT authentication with refresh tokens, role-based access control, audit logging, Raft consensus for state replication and Docker integration.
    \item \textbf{Chapter 4. Testing and Experimental Results}: presents the testing methodology including automated unit tests for the load balancer, scheduler and access control components, as well as manual end-to-end test scenarios covering node discovery, service deployment, fault tolerance, security enforcement and audit logging.
    \item \textbf{Chapter 5. Conclusions}: summarizes the achieved objectives, compares the platform with existing solutions, discusses its limitations and proposes future improvements such as persistent storage, service networking and auto-scaling.
\end{itemize}